\documentclass{article}
\usepackage{graphicx} % Required for inserting images
\usepackage{amsmath, amssymb, amsthm}
\usepackage{hyperref}

\usepackage{color}


\begin{document}




\hypersetup{
  colorlinks=true,
  linkcolor=blue,
  citecolor=blue,
  urlcolor=blue
}



\title{Honest Confidence Intervals over an \(\ell^1\) Slope Class in RKHS Linear Functional Regression}
\author{}
\date{}

\newcommand{\R}{\mathbb{R}}
\newcommand{\E}{\mathbb{E}}
\newcommand{\Prb}{\mathbb{P}}
\newcommand{\Supp}{\operatorname{Supp}}
\newcommand{\leb}{\operatorname{Leb}}
\newcommand{\supp}{\mathrm{supp}}

\maketitle







\newcommand{\cP}{\mathcal P}
\newcommand{\cX}{\mathcal X}
\newcommand{\cI}{\mathcal I}
\newcommand{\Leb}{\operatorname{Leb}}





\section{Statistical Experiment}

Let \(H=L^2[0,1]\), with inner product \(\langle\cdot,\cdot\rangle\) and norm
\(\|\cdot\|\). Let
\[
  \Theta_X=\{x\in H:0<\|x\|\le 1\},
  \qquad
  \Theta=\{\theta\in H:\|\theta\|\le 1\}.
\]
Fix a known compact, self-adjoint, nonnegative, and \emph{injective}
operator \(K:H\to H\), chosen so that the admissible class
\(\mathfrak P\) defined below is nonempty. Write \(B=K^{1/2}\). For \(x\in H\), define
\[
  z(x)=Bx.
\]
If \(X\sim P_X\) is supported on \(\Theta_X\), let \(P_Z=P_X\circ z^{-1}\)
be the induced law of \(Z=z(X)\), and define
\[
  T=T(P_X)=\E_{P_X}[Z\otimes Z].
\]

For each \(n\), the data are
\[
  D_n=\{(X_i,Y_i)\}_{i=1}^n,
  \qquad
  Y_i=\langle z(X_i),\theta_0\rangle+\epsilon_i,
\]
where \(X_i\stackrel{i.i.d.}{\sim}P_X\), \(\theta_0\in\Theta\), and
\(\epsilon_i\stackrel{i.i.d.}{\sim}N(0,1)\) are independent of the covariates.

An interval procedure is a measurable map
\[
  \varphi_n:\Theta_X\times (H\times\R)^n\to\cI,
\]
where \(\cI\) is identified with
\[
  \{(\ell,u)\in\R^2:\ell\le u\}
\]
equipped with the Borel sigma-field. The interval represented by
\((\ell,u)\) is \([\ell,u]\). The procedure may use \(D_n\), the prediction
point \(x\), the known operator \(K\), the confidence level \(\alpha\), and
the known constants in the assumptions. It may not use \(P_X\), \(P_Z\),
\(T\), or population eigenvalues/eigenfunctions, except through empirical
quantities computed from \(D_n\).

\section{Pointwise Classes, Global Risk, and Honesty}

Fix \(r>1/2\), \(\kappa\in(0,1)\), and a prediction point
\(x\in\Theta_X\). Write \(z=z(x)\). Let
\[
  \cP_n(z)=\cP_n(r,\kappa,z)
\]
be an increasing sequence of sets of laws \(P_X\) on \(\Theta_X\):
\(\cP_n(z)\subseteq\cP_{n+1}(z)\). The constants below do not depend on
\(n\), \(P_X\), \(j\), \(z\), or \(\theta\).

\begin{enumerate}
  \item[(A1)] There is \(C_4<\infty\) such that, for every
  \(P_X\in\cP_n(z)\) and every \(\theta\in\Theta\),
  \[
    \E\big[\langle Z,\theta\rangle^4\big]
    \le
    C_4\{\E[\langle Z,\theta\rangle^2]\}^2 .
  \]

  \item[(A2)] For every \(P_X\in\cP_n(z)\), the nonzero eigenvalues of
  \(T(P_X)\), written in decreasing order as
  \(s_1>s_2>\cdots>0\), are simple and have normalized eigenfunctions
  \(\psi_j\). There are constants \(0<a_s\le A_s<\infty\) such that
  \[
    a_s j^{-2r}\le s_j\le A_s j^{-2r}
    \qquad\text{for all }j\ge 1.
  \]

  \item[(A3)] There is \(C_\infty<\infty\) such that, for every
  \(P_X\in\cP_n(z)\) and every \(j\ge 1\),
  \[
    |\langle z,\psi_j\rangle|
    \le C_\infty n^\kappa\sqrt{s_j}.
  \]
  In addition, the same \(P_X\) satisfies the global random-covariate bound
  \[
    |\langle Z,\psi_j\rangle|
    \le C_\infty n^\kappa\sqrt{s_j}
    \qquad\text{for \(P_X\)-almost every \(X\) and all }j\ge 1.
  \]
  The first part of (A3) is the fixed-\(z\) condition used for pointwise
  coverage. The second part is the global \(P_X\)-condition used to control
  integrated squared length.

  \item[(A4)] (\textbf{Injectivity of \(T\).}) For every
  \(P_X\in\cP_n(z)\),
  \[
    \ker T(P_X)=\{0\}.
  \]
  Equivalently, the eigenfunctions \(\{\psi_j\}_{j\ge1}\) of \(T(P_X)\)
  form an orthonormal basis of \(H\), so \((\ker T(P_X))^{\perp}=H\)
  and every \(z(x)\) automatically lies in
  \(\overline{\operatorname{span}}\{\psi_j:j\ge1\}\). This excludes
  finite-rank operators \(T(P_X)\) and rules out unidentifiable slope
  components in \(\ker T(P_X)\).
\end{enumerate}

Define the global risk class
\[
  \cP_n=\bigcup_{x\in\Theta_X}\cP_n(z(x)).
\]
The squared-length risk is
\[
  R_n(\varphi_n;P_X,\theta_0)
  =
  \E_{\theta_0,P_X}
  \E_{X\sim P_X}
  \left[
    \Leb\{\varphi_n(X;D_n)\}^2
  \right],
\]
where the outer expectation is over the training sample \(D_n\), and the
inner expectation is over an independent test point \(X\sim P_X\).

A procedure has squared-length rate \(n^{-\rho(r)}\) on a slope class
\(\Theta_0(P_X)\) if
\[
  \limsup_{n\to\infty}
  \sup_{P_X\in\cP_n}
  \sup_{\theta_0\in\Theta_0(P_X)}
  n^{\rho(r)}R_n(\varphi_n;P_X,\theta_0)<\infty .
\]
An arbitrarily small exponent loss is considered matching: an upper bound
\(O(n^{-\rho(r)+\delta})\) for every \(\delta>0\), together with a lower bound
\(\Omega(n^{-\rho(r)})\), counts as the same minimax exponent.

For fixed \(x\in\Theta_X\), honesty means
for some given type-I error $\alpha\in(0,1)$,
\[
  \liminf_{n\to\infty}
  \inf_{P_X\in\cP_n(z(x))}
  \inf_{\theta_0\in\Theta_0(P_X)}
  \Prb_{\theta_0,P_X}
  \left[
    \langle z(x),\theta_0\rangle\in\varphi_n(x;D_n)
  \right]
  \ge 1-\alpha .
\]
The point \(x\) is fixed outside the probability statement. This is not an
averaged coverage statement over a fresh test point.

If \(\cP_n(z(x))=\varnothing\) eventually, the honesty condition at that
\(x\) is interpreted as vacuous. Equivalently, nonvacuous honesty is required
only for prediction points \(x\) for which \(\cP_n(z(x))\neq\varnothing\)
along an infinite tail or along the adversarial sequence under consideration.

All \(O(\cdot)\), \(o(\cdot)\), \(\Omega(\cdot)\), \(\lesssim\), and
\(\gtrsim\) constants must be uniform over the classes and prediction points
specified in each question.

\subsection*{Admissible Procedures and the Minimax Rate}

An interval procedure \(\varphi_n\) is \emph{admissible at level
\(\alpha\)} if it satisfies the honesty condition above
\emph{simultaneously} for every admissible problem instance and every
nonvacuous fixed prediction point \(x\in\Theta_X\) in it. Formally,
\[
  \mathcal A_\alpha
  :=
  \left\{
    \varphi_n :\
    \forall\,(\{\cP_n(\cdot)\},\Theta_0)\in\mathfrak P,\
    \forall\,x\in\Theta_X\text{ with }\cP_n(z(x))\neq\varnothing
    \text{ eventually},\
    \varphi_n\text{ satisfies the honesty condition for }(x,\{\cP_n(\cdot)\},\Theta_0)
  \right\}.
\]
The honesty condition is taken to be the definition of admissibility:
\(\mathcal A_\alpha\) and the honest-procedure class coincide. We
assume \(\mathcal A_\alpha\neq\varnothing\); a trivial witness is the
constant-width interval \(\varphi_n(x;D_n)=[-C\|x\|,\,C\|x\|]\) with
\(C\geq\|B\|\), which covers every admissible target with probability
one and hence is honest at every level \(1-\alpha\).

A \emph{problem instance} is a pair \((\{\cP_n(\cdot)\}_{n\ge1},\Theta_0)\)
where \(\{\cP_n(\cdot)\}\) is an increasing sequence of pointwise
distribution-class families satisfying the applicable subset of
(A1)--(A4) (and, where the question imposes them, (C1)--(C2)), and
\(\Theta_0(P_X)\) is a slope class as specified in the question. Let
\(\mathfrak P\) denote the collection of all such admissible problem
instances. The dependence on \(n\) is absorbed into the sequence
\(\{\cP_n(\cdot)\}_{n\ge1}\) inside the instance itself, so
\(\mathfrak P\) is not \(n\)-indexed.

Henceforth a procedure is a \emph{sequence}
\(\varphi=(\varphi_n)_{n\ge1}\) of single-\(n\) maps, and
\(\mathcal A_\alpha\) is interpreted as a class of such sequences.

The \emph{minimax rate exponent} \(\rho(r)\) is defined directly as the
largest exponent for which one fixed procedure sequence attains the rate
uniformly across all admissible problem instances:
\[
  \exists\,\varphi=(\varphi_n)_{n\ge1}\in\mathcal A_\alpha:\quad
  \limsup_{n\to\infty}\
  \sup_{(\{\cP_n(\cdot)\},\Theta_0)\in\mathfrak P}\
  \sup_{P_X\in\cP_n}\
  \sup_{\theta_0\in\Theta_0(P_X)}\
  n^{\rho(r)}\,
  R_n(\varphi_n;P_X,\theta_0)
  <
  \infty.
\]
Equivalently, with the infimum taken over procedure \emph{sequences}
outside the \(\limsup\),
\[
  \inf_{\varphi\in\mathcal A_\alpha}\
  \limsup_{n\to\infty}\
  \sup_{(\{\cP_n(\cdot)\},\Theta_0)\in\mathfrak P}\
  \sup_{P_X\in\cP_n}\
  \sup_{\theta_0\in\Theta_0(P_X)}\
  n^{\rho(r)}\,
  R_n(\varphi_n;P_X,\theta_0)
  <
  \infty.
\]
Note the order: the inf over sequences is taken \emph{outside} the
\(\limsup\) and all suprema, so the same \(\varphi=(\varphi_n)_n\) must
work for every \(n\), every instance, and every \((P_X,\theta_0)\). This
is strictly stronger than \(\limsup_n\inf_{\varphi_n}\) (which would
allow \(\varphi_n\) to be retuned per \(n\) per instance) and reflects
the universal-procedure obligation below.

\paragraph{Convention on empty classes.} Throughout, suprema over an
empty index set are interpreted as \(0\) and infima over an empty
index set as \(+\infty\). In particular, if no admissible problem
instance exists eventually (\(\mathfrak P=\varnothing\)) the rate
obligation is vacuous; if for a given instance \(\cP_n(z(x))=\varnothing\)
eventually for every \(x\), the integrated risk reduces to its inner
component on the residual index set, and any pointwise-honesty
requirement at such \(x\) is vacuous.

An arbitrarily small exponent loss is considered matching: an upper
bound of the above form with \(\rho(r)\) replaced by \(\rho(r)-\delta\)
for every \(\delta>0\), combined with a lower bound
\(\Omega(n^{-\rho(r)})\) for some single admissible problem instance,
counts as the same minimax exponent.

\paragraph{Upper-bound obligation.} A constructive upper bound must
exhibit \emph{one} procedure \(\varphi_n\in\mathcal A_\alpha\),
computable from \(D_n\), \(x\), \(K\), \(\alpha\), and the assumption
constants
\(\,r,\kappa,C_4,a_s,A_s,C_\infty,a_g\,\)
(all of which are known to \(\varphi_n\); any subset not relevant to a
given question may be omitted), that attains the rate uniformly over
every admissible problem instance. Except where a specific question
relaxes this rule (see Question 3a below), the procedure may not use
\(P_X\), \(P_Z\), \(T(P_X)\), or population eigenvalues/eigenfunctions
except through empirical quantities computed from \(D_n\).

\paragraph{Lower-bound obligation.} A lower bound need only exhibit
\emph{one} admissible problem instance
\((\{\cP_n^\star(\cdot)\},\Theta_0^\star)\in\mathfrak P\) such that
for every admissible procedure sequence
\(\varphi=(\varphi_n)_{n\ge1}\in\mathcal A_\alpha\),
\[
  \liminf_{n\to\infty}\
  \sup_{P_X\in\cP_n^\star}\
  \sup_{\theta_0\in\Theta_0^\star(P_X)}\
  n^{\rho(r)}\,
  R_n(\varphi_n;P_X,\theta_0)
  >0.
\]

\section{Questions}

\paragraph{Acceptable obligations for Question 3 (parts a and b).}
For both parts of Question 3, a proof is considered acceptable if it
establishes the stated rate under the following relaxed quantifier
structure on the smoothness exponent \(r\) and the (A3) envelope
exponent \(\kappa\):
\[
  \exists\,r_0<10\ \text{such that}\ \forall\,r\geq r_0,\
  \exists\,\kappa=\kappa(r)\in(0,1)
\]
for which the upper-bound obligation in the minimax block holds when
the admissible class \(\mathfrak P\) is restricted to instances whose
pointwise classes \(\cP_n(\cdot)\) use the smoothness exponent \(r\) and
the envelope exponent \(\kappa(r)\) in (A3).

\textbf{Strict constraints on \(\kappa(r)\).} The chosen \(\kappa(r)\)
must be a \emph{strictly positive constant}, in the following precise
sense:
\begin{itemize}
  \item \(\kappa(r)\) is a function of \(r\) \emph{only}; it does not
    depend on \(n\), on the sample \(D_n\), on \(P_X\), on \(x\), on
    any auxiliary parameter, or on any other quantity that varies with
    \(n\). In particular, \(\kappa(r)\) is \emph{not} permitted to be a
    sequence \(\kappa_n(r)\) tending to zero with \(n\).
  \item \(\kappa(r)>0\) strictly. Taking \(\kappa=0\), or taking
    \(\kappa_n\to 0\) along any subsequence, is explicitly excluded as
    it greatly trivializes the problem.
\end{itemize}
Equivalently: for each \(r\geq r_0\), the prover commits in advance to
a single positive number \(\kappa(r)\), and the rate must then hold
under (A3) armed with that fixed \(\kappa(r)\) for all \(n\).

The universal quantifier over
\((r,\kappa)\in(1/2,\infty)\times(0,1)\) in the original obligation is
thus replaced by the existential quantifier above on \(r_0\) and the
prover's chosen constant \(\kappa(r)\), with \(r\geq r_0\).

This relaxation does \emph{not} apply to Question 4: the strengthening
from \(\liminf\) to exact \(=1-\alpha\) is unrelated.

\paragraph{Question 3 --- shared setup for parts (a) and (b).}
Do \emph{not} impose any analogue of a nondegeneracy condition.
Instead, both parts assume \emph{(A1)--(A4)} (the base assumptions from
\textsection2) \emph{together with} the following two additional
conditions:

\begin{enumerate}
  \item[(C1)] There is \(a_g>0\) such that, for every \(P_X\in\cP_n\) and
  every \(j\ge 1\),
  \[
    s_j-s_{j+1}\ge a_g j^{-2r-1}.
  \]

  \item[(C2)] The slope class is
  \[
    \Theta_0(P_X)=\Theta'(P_X)
    :=
    \left\{
      \theta\in\Theta:
      \sum_{j=1}^{\infty}|\langle\theta,\psi_j\rangle|\le 1
    \right\}.
  \]
\end{enumerate}

The admissible class \(\mathfrak P^{(3)}\) for both parts of Question 3
is therefore the set of problem instances satisfying (A1)--(A4) and
(C1)--(C2), with \(\Theta_0(P_X)=\Theta'(P_X)\).

\paragraph{Procedure classes for Question 3.}
Let \(\mathcal A_\alpha^{(3b)}\) denote the subclass of \(\mathcal A_\alpha\)
consisting of procedure sequences with the original empirical
signature
\[
  \varphi_n:\Theta_X\times(H\times\R)^n\to\cI,
\]
honest at level \(\alpha\) for every instance in \(\mathfrak P^{(3)}\)
and every nonvacuous fixed prediction point in it.

For the oracle version, define a separate \emph{oracle} procedure
class \(\mathcal A_\alpha^{\mathrm{oracle}}\) whose elements are
measurable sequences
\[
  \varphi_n^{\mathrm{or}}:\Theta_X\times(H\times\R)^n\times\mathcal T\to\cI,
\]
where \(\mathcal T\) is the space of admissible population covariance
operators (compact, self-adjoint, nonnegative on \(H\), arising as
\(T(P_X)\) for some \(P_X\) in some instance of \(\mathfrak P^{(3)}\)).
Honesty for \(\mathcal A_\alpha^{\mathrm{oracle}}\) is the same as for
\(\mathcal A_\alpha\) except that the procedure is evaluated with
\(T(P_X)\) supplied as an additional argument; equivalently,
\(\varphi_n^{\mathrm{or}}(x;D_n;T(P_X))\) takes the place of
\(\varphi_n(x;D_n)\) inside the coverage probability.

\paragraph{Rate obligation (procedure may depend on \(\delta\)).}
For each part, a proof is considered acceptable if, for every
sufficiently small \(\delta>0\), there exists a procedure
\(\varphi^{(\delta)}\) in the relevant class
(\(\mathcal A_\alpha^{(3b)}\) for Question 3b,
\(\mathcal A_\alpha^{\mathrm{oracle}}\) for Question 3a) satisfying
\[
  \sup_{(\{\cP_n(\cdot)\},\Theta')\in\mathfrak P^{(3)}}\,
  \sup_{P_X\in\cP_n}\,\sup_{\theta_0\in\Theta'(P_X)}
  R_n(\varphi_n^{(\delta)};P_X,\theta_0)
  \;=\;
  O\!\left(n^{-2r/(2r+1)+\delta}\right).
\]
The procedure \(\varphi^{(\delta)}\) may depend on \(\delta\); it is
\emph{not} required that a single \(\delta\)-free procedure achieve
all such rates simultaneously.

In both parts the interval must be derived from an explicit limiting
distribution. The analysis must show that the bias divided by the
stochastic standard deviation is \(o_{\Prb}(1)\) \emph{pointwise} in
the fixed prediction point \(x\) and \emph{uniformly} over
\((P_X,\theta_0)\); precisely, for every nonvacuous fixed
\(x\in\Theta_X\),
\[
  \sup_{P_X\in\cP_n(z(x))}\,\sup_{\theta_0\in\Theta'(P_X)}\,
  \Prb_{\theta_0,P_X}\!\left[
    \frac{|b_n(x,P_X,\theta_0)|}{\sigma_n(x,P_X)}>\eta_n
  \right]\to0
\]
for some \(\eta_n\to0\), where \(b_n\) and \(\sigma_n\) are the
bias and stochastic standard deviation of the interval's underlying
estimator (defined as part of the construction).

The two parts differ only in which procedure class
\(\varphi^{(\delta)}\) lives in.

\paragraph{Question 3a (oracle version: procedure may use \(T\)).}
\(\varphi^{(\delta)}\in\mathcal A_\alpha^{\mathrm{oracle}}\). Concretely,
\(\varphi_n^{(\delta)}\) may depend on \(T(P_X)\) and on its eigenvalues
\(s_j\) and eigenfunctions \(\psi_j\), in addition to \(D_n\), \(x\),
\(K\), \(\alpha\), \(\delta\), and the assumption constants. The ``no
population objects'' clause of the upper-bound obligation in the
minimax block is therefore relaxed for Question 3a only, and replaced
by the explicit oracle signature above.

This is a strictly easier problem than Question 3b; it isolates the
population oracle / charging inequality from any empirical
spectral-transfer or eigenfunction-perturbation problem. A solution to
Question 3a should make clear which steps would require empirical
transfer in Question 3b and which would remain unchanged.

\paragraph{Question 3b (empirical version: original Question 3).}
\(\varphi^{(\delta)}\in\mathcal A_\alpha^{(3b)}\). Concretely,
\(\varphi_n^{(\delta)}\) is computable from \(D_n\), \(x\), \(K\),
\(\alpha\), \(\delta\), and the assumption constants only. In
particular, \(\varphi_n^{(\delta)}\) may \emph{not} depend on
\(T(P_X)\), \(P_X\), \(P_Z\), or on the population
eigenvalues/eigenfunctions \(s_j,\psi_j\), except through empirical
quantities computed from \(D_n\). No additional empirical perturbation
assumptions may be imposed: any needed sample-covariance, empirical
eigenvalue, or empirical eigenfunction approximation bounds must be
derived from (A1)--(A4), (C1)--(C2), and the sample.

This is the original (harder) form of the question. Solving Question
3a first is acceptable as a partial result.

\paragraph{Remark.}
The bias/SD requirement is to be read for nonvacuous fixed prediction
points \(x\in\Theta_X\), i.e.\ those for which \(\cP_n(z(x))\) is
eventually nonempty. (Note \(\Theta_X\) already excludes \(x=0\), so
\(z(x)\neq 0\) holds throughout under (A4) plus \(K\) injective.) For
such \(x\), under (A4) and Gaussian noise, the stochastic standard
deviation \(\sigma_n(x)\) of any reasonable estimator-based interval
is strictly positive at every finite \(n\) with probability one even
though \(\sigma_n(x)\to 0\); the \(o_{\Prb}(1)\) statement is the
standard ratio in this finite-\(n\) sense, where the bias is required
to vanish faster than the SD. No population-level positive-variance
assumption is imposed for either part of Question 3.

\paragraph{Question 4.}
The honesty definition above uses a \(\liminf\). In this question the
procedure \(\varphi_n\) is required to be \emph{empirical}:
\(\varphi_n\in\mathcal A_\alpha^{(3b)}\), i.e.\ computable from \(D_n\),
\(x\), \(K\), \(\alpha\), and the assumption constants only, with no
dependence on \(T(P_X)\) or population eigenvalues/eigenfunctions
except through empirical quantities. \(\varphi_n\) need not be the
same procedure that solved Question 3a or 3b.

Identify natural additional conditions, imposed only on
population-level objects \(P_X\), \(\Theta_0(P_X)\), and \(T(P_X)\),
under which this can be strengthened to
\[
  \lim_{n\to\infty}
  \inf_{P_X\in\cP_n(z(x))}
  \inf_{\theta_0\in\Theta_0(P_X)}
  \Prb_{\theta_0,P_X}
  \left[
    \langle z(x),\theta_0\rangle\in\varphi_n(x;D_n)
  \right]
  =
  1-\alpha
\]
for every \(x\in\Theta_X\) for which the pointwise class is nonvacuous.
The equality is required to be \emph{exact} in the uniform two-sided
sense:
\[
  \lim_{n\to\infty}\
  \sup_{P_X\in\cP_n(z(x))}\,\sup_{\theta_0\in\Theta_0(P_X)}
  \left|\,
    \Prb_{\theta_0,P_X}
    \!\left[\langle z(x),\theta_0\rangle\in\varphi_n(x;D_n)\right]
    -(1-\alpha)
  \,\right|
  \;=\;0,
\]
for every nonvacuous fixed prediction point \(x\). A procedure that
achieves only \(\liminf_n\inf_{\cdots}\Prb[\cdots]\ge 1-\alpha\),
allowing coverage strictly above \(1-\alpha\) on a non-negligible
subset of admissible instances, does not satisfy Question 4: the
interval must not be asymptotically conservative on any tail.

This $\varphi_n$ and its analysis can be different from previous questions.
The
additional conditions may not involve empirical quantities such as the sample
covariance operator. Explain why the proposed conditions are sufficient and
whether they are close to necessary.

\paragraph{Acceptable partial resolutions.}
The three sub-questions \emph{(3a)}, \emph{(3b)}, and \emph{(4)} are
treated as \emph{independent} acceptable targets. A complete answer
to \emph{any one} of them, with a rigorous proof, counts as a full
acceptable resolution of the problem; the prover need not solve all
three. Specifically:
\begin{itemize}
  \item Solving Question~3a alone (the oracle version with population
    \(T\) available to the procedure) is acceptable. Question~3a is
    deliberately the easiest of the three.
  \item Solving Question~3b alone (the empirical version) is
    acceptable.
  \item Solving Question~4 alone (the exact-coverage strengthening) is
    acceptable, treating the \(\varphi_n\) of Question~4 as an
    independent construction in the empirical procedure class
    \(\mathcal A_\alpha^{(3b)}\).
  \item A negative answer (rigorous lower bound or counterexample)
    for any one of \emph{(3a)}, \emph{(3b)}, \emph{(4)} is also
    acceptable as a full resolution of that part, with the same
    rules: a single part proved negative counts as a complete
    resolution.
\end{itemize}
The prover should choose the most tractable target and resolve it
cleanly rather than attempt all three simultaneously. The final
proof document must clearly state which sub-question is being
resolved (or that all three are) and which form (constructive
upper bound, lower bound / counterexample) applies.

\paragraph{Tractability hierarchy (advisory).}
On the basis of the original (paper) literature and on prior runs of
this pipeline, the following tractability ordering is expected:
\(\text{(3a)}\) is strictly easier than \(\text{(3b)}\) (the
empirical-eigenfunction transfer is exactly the hard part of
\(\text{(3b)}\)). Question~\(\text{(4)}\) is independent of
\(\text{(3a)}\)/\(\text{(3b)}\) and asks a coverage-calibration
question that may be addressed without solving the rate problem at
all. Question~\(\text{(3b)}\) in its fully empirical form is the
hardest of the three and is the one for which a counterexample
\emph{may} be the right answer (see the prior obstruction to an
\(o(n^{-1/2})\) empirical eigenfunction tail-leakage bound that
appears in standard analyses; the prover may either find a route
that avoids this obstruction or prove that no such route can exist).




\newpage




\end{document}

